\section{Segunda iteración de la Fase 5: Formato del Cartucho}
\label{sec:F5I2}

El análisis de la iteración anterior dejó el inventario cerrado: en
el GWC viaja todo lo que los dos motores leen en un \textit{tick}, y
nada más. Esta iteración define la estructura exacta del
fichero que se va a generar intentando optimizarlo al máximo.

\subsection{Datos numéricos y el valor nulo}
\label{sec:F5I2Nulo}

Los enteros viajan en \textbf{little-endian} (menos significativo
primero), el orden natural del procesador del ESP32: el firmware lee
los bytes en una estructura de C sin conversión.\footnote{Aunque se debe tener en cuenta que en un ESP32, en \textit{offsets} impares un \textit{cast} a \texttt{uint16\_t*} puede dar excepción de alineación.}

Usaremos un solo convenio para todo el fichero:

\begin{itemize}
	\item Todos los punteros e índices del formato son enteros sin signo, que comienzan
	por el valor \textbf{cero}.
	
	\item Usaremos el valor \textbf{nulo} para indicar que el enlace o el elemento
	\emph{no existen} (un segmento sin siguiente, un sprite ausente).
\end{itemize}

La representación del valor \textbf{nulo} depende de la anchura del campo, pero siempre tendrá todos sus bits a 1. Como consecuencia, un campo de 1 byte direcciona hasta 255 elementos (0 a 254, el nulo queda reservado).

\begin{table}[H]
	\centering
	\begin{tabular}{|c|c|}
		\hline
		\textbf{Anchura del campo} & \textbf{Valor nulo} \\
		\hline
		1 byte  & \texttt{0xFF} \\ \hline
		2 bytes & \texttt{0xFFFF} \\ \hline
		4 bytes & \texttt{0xFFFFFFFF} \\ \hline
	\end{tabular}
	\caption{Representación del valor nulo.}
\end{table}

\subsection{Estructura del fichero}
\label{sec:F5I2Estructura}

El fichero se estructura en un registro de cabecera que lo identifica, seguido de varios bloques de datos de longitud variable, algunos opcionales. Todos los bloques usan la misma estructura:

\begin{table}[H]
	\centering
	\begin{tabular}{|l|c|c|p{8.4cm}|}
		\hline
		\textbf{Campo} & \textbf{Tipo} & \textbf{Bytes} & \textbf{Contenido} \\
		\hline
		
		Tipo del bloque & Num & 1 & Un número que define el tipo:
		\begin{itemize}
			\item 0 Datos Generales
			\item 1 Configuración
			\item 2 Sprites
			\item 3 Segmentos
			\item 4 Secuencias (opcional)
			\item 5 Fondo (opcional)
		\end{itemize} \\ \hline
		
		Tamaño & Num & 4 & Tamaño \underline{del contenido} del bloque (\textbf{sin incluir} los 5 bytes de cabecera).\\ \hline
		
		Contenido & Variable & Variable & Contenido del bloque \\
		
		\hline\hline
		\multicolumn{4}{c}{\textbf{Tamaño: Variable}} \\ \hline
	\end{tabular}
	\caption{Definición de Bloques.}
	\label{tab:definicionBloques}
\end{table}

\subsection{La cabecera}
\label{sec:F5I2Cabecera}

El fichero comienza por una cabecera de tamaño fijo que lo identifica unívocamente:

\begin{table}[H]
	\centering
	\begin{tabular}{|l|c|c|p{8.4cm}|}
		\hline
		\textbf{Campo} & \textbf{Tipo} & \textbf{Bytes} & \textbf{Contenido} \\
		\hline
		
		Identificador & Char & 3 & Los caracteres «GWC» como identificación \\ \hline
		
		Versión & Char & 1 & Versión del formato, inicialmente «1» \\ \hline
		
		CheckSum & Char & 6 & Suma de verificación de 16 bits del resto del fichero tras esta cabecera, sin incluirla, en formato «\#hhhh\#» para que sea visible\footnotemark \\ \hline
		
		Tamaño & Num & 4 & Tamaño del fichero \\ \hline	
		\hline\hline
		\multicolumn{4}{c}{\textbf{Tamaño: 14 Bytes}} \\ \hline
	\end{tabular}
	\caption{Campos de la cabecera del GWC.}
	\label{tab:cabeceraGWC}
\end{table}
\footnotetext{Se incluye en la generación del GWI el campo «CódigoGWC» que incluye este mismo valor como referencia de que ambos se han generado desde el mismo código.}


\subsection{Bloque 0: Datos Generales}
\label{sec:F5I2DatosGenerales}

Este bloque define los datos generales del cartucho y el juego, siendo una estructura de tamaño fijo:

\begin{table}[htbp]
	\centering
	\begin{tabular}{|l|c|c|p{8.4cm}|}
		\hline
		\textbf{Campo} & \textbf{Tipo} & \textbf{Bytes} & \textbf{Contenido} \\
		\hline
		Tipo del bloque & Num & 1 & Número fijo: 0\\
		\hline
		
		Tamaño & Num & 4 & Tamaño del bloque. 35 fijo \\
		\hline
		
		Contenido & Variable & 35 & Contenido del bloque \\
		
		\hline\hline
		\multicolumn{4}{c}{\textbf{Tamaño: 40 Bytes}} \\ \hline
	\end{tabular}
	\caption{Bloque de Datos Generales.}
\end{table}

La estructura de este bloque está pensada para que el ESP32 la lea entera en una
estructura de C y sepa dónde está todo sin recorrer nada (se indica expresamente los bloques obligatorios y los opcionales):

\begin{table}[H]
	\centering
	\begin{tabular}{|l|c|c|p{7.4cm}|c|}
		\hline
		\textbf{Campo} & \textbf{Tipo} & \textbf{Bytes} & \textbf{Contenido} & \textbf{Bloque} \\
		\hline
		
		Estilo & Num & 1 & Estilo necesario: Se comprueba que el cartucho es adecuado para esa máquina & -  \\ \hline
		Ancho & Num & 2 & Ancho de la pantalla al que el generador ha escalado el cartucho & -  \\ \hline
		Alto & Num & 2 & Igual para el alto de la pantalla & -  \\ \hline
		Sprites & Num & 1 & Número de sprites en el cartucho & -  \\ \hline
		Segmentos & Num & 1 & Número de segmentos en el cartucho & -  \\ \hline
		Secuencias & Num & 1 & Número de secuencias en el cartucho & - \\ \hline
		Fondo & Num & 3 & Color de fondo en RGB. & - \\ \hline
		Inicio Bloque 0 & Num & 4 & Inicio del bloque 0 & Ob. \\ \hline
		Inicio Bloque 1 & Num & 4 & Inicio del bloque 1 & Ob. \\ \hline
		Inicio Bloque 2 & Num & 4 & Inicio del bloque 2 & Ob. \\ \hline
		Inicio Bloque 3 & Num & 4 & Inicio del bloque 3 & Ob. \\ \hline
		Inicio Bloque 4 & Num & 4 & Inicio del bloque 4 o \textbf{nulo} & Opc \\ \hline
		Inicio Bloque 5 & Num & 4 & Inicio del bloque 5 o \textbf{nulo} & Opc \\
		
		\hline\hline
		\multicolumn{5}{c}{\textbf{Tamaño: 35 Bytes}} \\ \hline
	\end{tabular}
	\caption{Bloque 0: Datos Generales.}
	\label{tab:bloque0DatosGenerales}
\end{table}

Los punteros al inicio de los bloques se establecen solo para que sea
más rápida la lectura del cartucho e ir directos al inicio de cada uno de los bloques.

Si no hay secuencias o imagen de fondo, los punteros a dichos bloques tendrán el valor \textbf{nulo} como puntero. El campo \textbf{Fondo} se usará solo en caso de que no tengamos una imagen de fondo: este es el color que se usará como fondo de la imagen. Por defecto será el blanco.


\subsection{Bloque 1: Configuración general del juego}
\label{sec:F5I2ConfGen}

\textbf{NOTA:} Inicialmente eran 8 valores, pero se amplían a 9 al detectar un error, ver \ref{sec:F5I13}

Es necesario disponer de varios datos generales sobre el juego, velocidades, aceleraciones, lanzamientos, etc. Para guardar estos datos usaremos una tabla de \textbf{nueve} valores numéricos de 16 bits, de esta manera se reservan estos 18 bytes que estarán siempre presentes, y serán leídos por el firmware en una sola pasada como un array de \texttt{uint16\_t}.

\begin{table}[htbp]
	\centering
	\begin{tabular}{|l|c|c|p{8.4cm}|}
		\hline
		\textbf{Campo} & \textbf{Tipo} & \textbf{Bytes} & \textbf{Contenido} \\
		\hline
		Tipo del bloque & Num & 1 & Número fijo: 1\\
		\hline
		
		Tamaño & Num & 4 & Tamaño del bloque. \sout{16} 20 fijo \\
		\hline
		
		Contenido & Num & 18 & Contenido del bloque, 9 valores numéricos \\
		
		\hline\hline
		\multicolumn{4}{c}{\textbf{Tamaño: 23 Bytes}} \\ \hline
	\end{tabular}
	\caption{Bloque de Configuración general del juego.}
\end{table}

Cada campo contiene
el valor que se definió en el diseñador y, si no se hizo, el valor \sout{por
defecto ya definido en el GWI} nulo\footnote{ver \ref{sec:F5I13}} en los campos 0, 1 y 5 de la lista de posibles valores. De esta manera el campo viaja siempre, sea cual
sea su origen:

\begin{table}[H]
	\centering
	\begin{tabular}{|c|l|p{9.5cm}|}
		\hline
		\textbf{Nro} & \textbf{Contenido} & \textbf{Descripción} \\
		\hline
		
		00 & Pausa del simulador & La pausa por plataforma, en milisegundos, para el navegador; \textbf{nulo} si no aplica  \\ \hline
		
		01 & Pausa del externo & La misma pausa para el firmware del ESP32: mismos datos, distinta puerta; \textbf{nulo} si no aplica  \\ \hline
		
		02 & Acelerar cada & Cada cuántos ticks se produce el aumento de velocidad\\ \hline
		
		03 & Resta & Milisegundos que se quitan en cada aumento\\ \hline
		
		04 & Mínimo & Pausa a la que la aceleración se detiene\\ \hline
		
		05 & Aumentar cada & Número de lanzamientos entre aumentos de la política de rutas; \textbf{nulo} si no aplica  \\ \hline
		
		06 & Mecanismo de fin & Por contador o por familia de iconos \\ \hline
		
		07 & Contador de fin & El contador elegido cuando el mecanismo es por contador; \textbf{nulo} si no aplica \\ \hline
		
		08 & Reglas del juego & Campo de bits con las reglas del cartucho: el bit 0 es la regla \textbf{combinar} (dos pulsaciones en un mismo tick)\footnote{Inicialmnete no se maneja esta regla pues será una posible ampliación}; los demás bits están reservados y valen 0. Vale 0 si el cartucho no declara ninguna regla \\ \hline
			
	\end{tabular}
	\caption{Datos en la tabla de configuración general del juego.}
	\label{tab:configuracionGeneral}
\end{table}


\subsection{Bloque de Sprites}
\label{sec:F5I2Sprites}

Los sprites del GWI llegan ya deduplicados y en blanco y negro
puro; el generador optimiza la ocupación convirtiendo los PNG en máscaras de 1 bit escaladas a
la pantalla declarada. El bloque tiene esta estructura:

\begin{table}[htbp]
	\centering
	\begin{tabular}{|l|c|c|p{8.4cm}|}
		\hline
		\textbf{Campo} & \textbf{Tipo} & \textbf{Bytes} & \textbf{Contenido} \\
		\hline
		Tipo del bloque & Num & 1 & Número fijo: 2 \\ \hline
		Tamaño & Num & 4 & Tamaño del bloque. \\ \hline
		Contenido & Num & Variable & Contenido del bloque \\
		
		\hline\hline
		\multicolumn{4}{c}{\textbf{Tamaño: Variable}} \\ \hline
	\end{tabular}
	\caption{Bloque de Sprites.}
\end{table}

% evita corte de tabla
\Needspace{18\baselineskip}
Dentro del bloque de sprites los datos se almacenan en dos zonas: la tabla de direcciones, con la posición de cada sprite, y las máscaras, con el contenido de cada uno. Cada dirección señala el inicio de su máscara:

\begin{table}[H]
	\centering
	\begin{tabular}{|l|c|c|p{8.4cm}|}
		\hline
		\textbf{Campo} & \textbf{Tipo} & \textbf{Bytes} & \textbf{Contenido} \\
		\hline
		Dirección & Num & 4 & Posición de la máscara de un sprite dentro
		del fichero; una entrada por sprite, en el orden de la tabla del
		GWI \\
		\hline
		\multicolumn{4}{|c|}{\textit{La máscara señalada por cada
				dirección:}} \\
		\hline
		Ancho & Num & 2 & Ancho del sprite en píxeles, ya escalado \\
		\hline
		Alto & Num & 2 & Alto del sprite en píxeles, ya escalado \\
		\hline
		Fila & Num & $n \times h$ & Las $h$ filas del sprite ($h$ = alto),
		de arriba abajo, con $n = \lceil \text{ancho}/8 \rceil$ bytes por
		fila: 8 píxeles por byte, el bit más significativo a la izquierda,
		blanco $= 0$, negro $= 1$ \\
		\hline
	\end{tabular}
	\caption{La tabla de sprites y la estructura de cada máscara.}
\end{table}

Un sprite de $16 \times 10$ ocupa $n = 2$ bytes por fila y
$2 \times 10 = 20$ bytes de máscara, más sus dimensiones: los PNG
de cientos de bytes han desaparecido. Las máscaras salen del mismo
generador de huellas que usa la deduplicación de la autoría, así
que los sprites siguen compartidos: los segmentos que repiten el
mismo dibujo apuntan al mismo índice de la tabla, y la ocupación es
la de los sprites \emph{distintos}, no la de las apariciones.

\subsection{Bloque de segmentos}
\label{sec:F5I2Blsegmentos}

Lo que define el mecanismo del juego, por tanto es el corazón del formato. Se ha optado por \textbf{limitar su número a 255}, para que los enlaces en la tabla ocupen solo 1 byte. Esto es suficiente para los juegos que se han analizado, aunque en algún juego de doble pantalla como el Donkey Kong puede estar muy al límite, por tanto \textbf{el generador dará un error bloqueante si hay más de 255 segmentos}; si fuera necesario ampliarlo esto se puede ajustar en una nueva versión del fichero, que aumente los contadores de segmentos a 2 bytes.

La estructura general de este bloque es:

\begin{table}[H]
	\centering
	\begin{tabular}{|l|c|c|p{8.4cm}|}
		\hline
		\textbf{Campo} & \textbf{Tipo} & \textbf{Bytes} & \textbf{Contenido} \\
		\hline
		
		Tipo del bloque & Num & 1 & Número fijo: 3\\ \hline
		
		Tamaño & Num & 4 & El tamaño total del bloque. \\ \hline
		
		Contenido & Num & $32 \times n$ & Contenido del bloque, n registros de 32 bytes\\
		
		\hline\hline
		\multicolumn{4}{c}{\textbf{Tamaño: Variable}} \\ \hline
	\end{tabular}
	\caption{Bloque de segmentos.}
	\label{tab:bloqueSegmentos}
\end{table}

Cada segmento es un \textbf{registro de tamaño fijo} de 32 bytes: la posición del segmento
$i$ es \emph{dirección del segmento i} = $\text{inicio contenido} + (i \times \text{tamaño del registro})$, donde en esta versión $\text{tamaño del registro} = 32$, sin
recorrer nada — la misma lectura directa que hace hoy el motor con el array
de JavaScript. Los campos son exactamente los que el motor de estados y el
motor gráfico leen en un tick, en la forma en que los leen:

\begin{table}[H]
	\centering
	\begin{tabular}{|l|c|c|p{8.4cm}|}
		\hline
		\textbf{Campo} & \textbf{Tipo} & \textbf{Bytes} & \textbf{Contenido} \\
		\hline
		
		Tipo & Num & 1 & El tipo del segmento (jugador, juego, inicio, vida,
		muerte, dígito, contador, adicional, inicio adicional): la misma enumeración con nombre que vive en el puente \\
		\hline
		Posición x & Num & 2 & Posición del segmento en la imagen del fondo escalada a la pantalla \\
		\hline
		Posición y & Num & 2 & Posición del segmento en la imagen del fondo escalada a la pantalla \\
		\hline
		Sprite & Num & 1 & Índice en la tabla de sprites \\
		\hline
		\emph{siguiente} & Num & 1 & Enlace de avance de rutas y cadenas \\
		\hline
		\emph{salvado} & Num & 1 & Cadena del desenlace con jugador \\
		\hline
		\emph{muerte} & Num & 1 & Cadena del desenlace sin jugador \\
		\hline
		\emph{salvadoContador} & Num & 1 & El contador que se ajusta cuando el desenlace es salvado (1..4); \textbf{nulo} si no se toca ninguno \\
		\hline
		\emph{muerteContador} & Num & 1 & El contador que se ajusta cuando el desenlace es muerte (1..4); \textbf{nulo} si no se toca ninguno \\
		\hline
		\emph{mostrarJunto} & Num & 1 & Enlace a su acompañante («Mostrar junto a») \\
		\hline
		\emph{izquierda, derecha, arriba, abajo} & Num & 4 & Las cuatro direcciones del jugador (1 byte cada una) \\
		\hline
		\emph{colision} & Num & 1 & El segmento jugador con el que colisiona \\
		\hline
		\emph{barra\_a} \ldots{} \emph{barra\_g} & Num & 7 & Las siete barras de un dígito (1 byte cada una) \\
		\hline
		\emph{digito\_1} \ldots{} \emph{digito\_4} & Num & 4 & Los dígitos de un contador, \emph{digito\_1} = las unidades (1 byte cada una) \\
		\hline
		Propiedades & Num & 1 & Empaquetadas en bits:
		\begin{itemize}
			\item \textbf{Bit 0:} Primer jugador
			\item \textbf{Bit 1:} Cuenta ascendente
			\item \textbf{Bit 2:} Presentación del cero
			\item \textbf{Bits 3 a 7:} Libres a Cero
		\end{itemize} \\
		\hline
		Contador interno & Num & 1 & Número de contador (1--4) en los contadores; en el resto 0 \\
		\hline
		Valor inicial & Num & 1 & Valor de arranque del contador o Cero \\
		\hline
		Ruta & Num & 1 & Número de ruta en los inicios \\
		\hline
		\hline\hline
		\multicolumn{4}{|c|}{\textbf{Tamaño: 32 bytes}} \\
		\hline
	\end{tabular}
	\caption{El registro de segmento.}
\end{table}

Los enlaces usan el convenio definido para el valor \textbf{nulo}, por lo que el enlace no existe: un segmento puede no tener un enlace a siguiente o un enlace a derecha.

Con el Parachute (unos 40 segmentos) el bloque entero ocupa unos
1,3 KB. Lo que \textbf{no} hay en el registro son datos no necesarios: los identificadores legibles, los fantasmas tipo
99, los segmentos absorbidos sin visibilidad\footnote{Los segmentos que forman parte de un dígito no son visibles por defecto, pero sí deben viajar en el cartucho}, los enlaces inversos que el motor no lee
(el inverso calculado en \emph{mostradoJunto}, \emph{colisionDe}) y toda la sección de autoría. El generador los descarta al vuelo: un segmento que nadie
enciende ni pinta no existe en el GWC.

El escalado del que habla el análisis se aplica aquí: las
posiciones y los sprites se transforman con el muestreo por vecino
más cercano y escala por eje del simulador, y lo que viaja es el
resultado final en el espacio de la pantalla. El ESP32 recibe
números listos para volcar.

\subsection{Bloque de Secuencias}
\label{sec:F5I2Secuencias}

Puede haber varias \textbf{secuencias} en el juego, cada una con 5 segmentos posibles, tal como se definió en el GWI. Para añadir esto usaremos un bloque con registros de longitud fija de 6 bytes, con estos datos:

\begin{table}[H]
	\centering
	\begin{tabular}{|l|c|c|p{8.4cm}|}
		\hline
		\textbf{Campo} & \textbf{Tipo} & \textbf{Bytes} & \textbf{Contenido} \\
		\hline
		
		Tipo del bloque & Num & 1 & Número fijo: 4\\ \hline
		
		Tamaño & Num & 4 & El tamaño total del bloque. \\ \hline
		
		Contenido & Num & $6 \times n$ & Contenido del bloque, n registros\\
		
		\hline\hline
		\multicolumn{4}{c}{\textbf{Tamaño: Variable}} \\ \hline
	\end{tabular}
	\caption{Bloque de Secuencias.}
	\label{tab:bloqueSecuencias}
\end{table}

Los registros llevan sus \textit{ticks} de espera y
hasta cinco índices de segmento (con \textbf{nulo} cuando el
hueco está libre), lanzados en orden, uno por tick.

\begin{table}[H]
	\centering
	\begin{tabular}{|l|c|c|p{8.4cm}|}
		\hline
		\textbf{Campo} & \textbf{Tipo} & \textbf{Bytes} & \textbf{Contenido} \\
		\hline
		
		Ticks & Num & 1 & Cuántos Ticks hay entre secuencias \\ \hline
		Segmento 1 & Num & 1 & Qué segmento es el primero de la secuencia \\ \hline
		Segmento 2 & Num & 1 & Qué segmento es el segundo de la secuencia \\ \hline
		Segmento 3 & Num & 1 & Qué segmento es el tercero de la secuencia \\ \hline
		Segmento 4 & Num & 1 & Qué segmento es el cuarto de la secuencia \\ \hline
		Segmento 5 & Num & 1 & Qué segmento es el quinto de la secuencia \\
		
		\hline\hline
		\multicolumn{4}{c}{\textbf{Tamaño: 6 bytes}} \\ \hline
	\end{tabular}
	\caption{El registro de secuencias.}
	\label{tab:registroSecuencias}
\end{table}

Los enlaces usan el convenio definido: \textbf{nulo} = ausente.

\subsection{Bloque de la imagen de Fondo}
\label{sec:F5I2Fondo}

La imagen de fondo no existe en todas las máquinas, por ejemplo en Ball no existe, y aunque tampoco es imprescindible ya que es más un adorno, si queremos replicar un juego al máximo será conveniente incluirla. Para añadir la imagen de fondo usaremos un bloque con estos datos:

\begin{table}[H]
	\centering
	\begin{tabular}{|l|c|c|p{8.4cm}|}
		\hline
		\textbf{Campo} & \textbf{Tipo} & \textbf{Bytes} & \textbf{Contenido} \\
		\hline
		
		Tipo del bloque & Num & 1 & Número fijo: 5 \\ \hline
		
		Tamaño & Num & 4 & El tamaño total del bloque. \\ \hline
		
		Contenido & Variable & Variable & La imagen completa de la pantalla \\
		
		\hline\hline
		\multicolumn{4}{c}{\textbf{Tamaño: Variable}} \\ \hline
	\end{tabular}
	\caption{Bloque de la imagen del fondo.}
\end{table}

Si el bloque existe se usará como imagen de fondo, si no existe el fondo será una pantalla con el color indicado en el campo \textbf{Fondo} del Bloque 0, que es el blanco por defecto.

El formato de la imagen inicialmente usará el formato \textbf{RGB565} directo: 2 bytes por píxel en little-endian, $\text{Ancho} \times \text{Alto}$ píxeles según lo indicado en el Bloque 0, sin compresión ni decodificador — el formato nativo del TFT del ESP32, listo para volcar.

Pero reservaremos una iteración solo para buscar la mejor técnica para poder optimizar la ocupación de las imágenes, como usar compresión RLE, usar paletas de colores separadas de la imagen, o ya que la mayor parte de la pantalla no tiene nada en el fondo, ya que taparía a los segmentos, es muy prometedor usar Skip RLE, con un color de fondo predominante y guardar solo las áreas con otros colores.

\subsection{Generación determinista del cartucho}
\label{sec:F5I2Def}

La generación es siempre \textbf{determinista}: el mismo GWI debe
producir siempre el mismo GWC byte a byte, con la misma suma de
verificación. El generador es equivalente a un compilador, con
el mismo origen siempre produce el mismo resultado.

Para asegurar esto, se añade en el Cartucho GWI un campo denominado \textbf{codigoGWC} con el contenido del campo Checksum del cartucho GWC junto a sus marcas de inicio y fin (formato \#FFFF\#), que coincide cuando ambos cartuchos se han generado a partir de los mismos datos exactamente.

\subsection{Posibles mejoras}
\label{sec:F5I2Mejoras}

Para reducir más el tamaño del cartucho, se pueden usar en algunos lugares ciertas mejoras, pero en general cambiarían espacio por tiempo de procesamiento, lo que no se ve necesario en este momento y no se abordarán en esta fase.

La principal es que si el cartucho tiene pocos segmentos se pueden reducir el tamaño de los punteros: se podría usar el campo con el contador de segmentos para decidir el tamaño real de los punteros, que seguramente se puede reducir a solo 2 bytes en muchos juegos.

Los registros de segmento usan campos que no se emplean en todos los tipos; cada registro reserva siempre espacio para 7 barras, 4 dígitos, 4 direcciones, etc. Usando registros de longitud variable dependiendo del tipo se podrían llegar a reducir a segmentos de entre 14 bytes y 30 bytes, en lugar de usar siempre los 32 actuales, lo que puede llegar a representar hasta 1~KB menos en juegos con 80 segmentos.

